\section*{Conclusion and Further Work}
\addcontentsline{toc}{section}{Conclusion and Further Work}

An explicitly presented harmonic lattice gives a precise discrete-CVP comparison, while
the constructed common-input examples prevent uniform exact recovery from
the specified public data.


The comparison suggests studying encoded HCP families with explicit
distributions on hidden presentations and Hodge-structured noise. Their
computational difficulty must be assessed within that specified model.

Discrete Morse reduction for cellular sheaf cohomology
\cite{curryGhristNanda2015morse} suggests a computational refinement.
A chain-homotopy equivalence alone does not identify the original harmonic
representative or preserve its distance. An application to the original
recovery target needs explicit retained lifts and a metric comparison.

Further work includes extending the homological model to full first-order
fibrational semantics, realizing richer families topologically, and studying
transported observations.
